% DIE VERZAUBERUNG - Vollständiges Drehbuch-Paket
% Basierend auf dem Roman von Hermann Broch
% Drehbuch: Boris Marte

\documentclass[11pt,a4paper,twoside]{book}

% ============================================================
% PAKETE
% ============================================================

\usepackage[utf8]{inputenc}
\usepackage[T1]{fontenc}
\usepackage[ngerman]{babel}
\usepackage{geometry}
\usepackage{setspace}
\usepackage{titlesec}
\usepackage{fancyhdr}
\usepackage{graphicx}
\usepackage{xcolor}
\usepackage{booktabs}
\usepackage{longtable}
\usepackage{array}
\usepackage{tabularx}
\usepackage{enumitem}
\usepackage{tcolorbox}
\usepackage{fontspec}
\usepackage{microtype}
\usepackage{hyperref}
\usepackage{titletoc}
\usepackage{epigraph}

% ============================================================
% SEITENGEOMETRIE
% ============================================================

\geometry{
    a4paper,
    left=30mm,
    right=25mm,
    top=30mm,
    bottom=30mm,
    headheight=14pt
}

% ============================================================
% FARBEN
% ============================================================

\definecolor{darkgray}{RGB}{50,50,50}
\definecolor{midgray}{RGB}{100,100,100}
\definecolor{lightgray}{RGB}{220,220,220}
\definecolor{accentcolor}{RGB}{120,80,40}
\definecolor{boxbg}{RGB}{248,246,242}

% ============================================================
% SCHRIFTEN
% ============================================================

\setmainfont{TeX Gyre Termes}
\setsansfont{TeX Gyre Heros}

% ============================================================
% KOPF- UND FUSSZEILEN
% ============================================================

\pagestyle{fancy}
\fancyhf{}
\fancyhead[LE]{\small\textsc{Die Verzauberung}}
\fancyhead[RO]{\small\textsc{\leftmark}}
\fancyfoot[C]{\thepage}
\renewcommand{\headrulewidth}{0.4pt}

\fancypagestyle{plain}{
    \fancyhf{}
    \fancyfoot[C]{\thepage}
    \renewcommand{\headrulewidth}{0pt}
}

% ============================================================
% KAPITELÜBERSCHRIFTEN
% ============================================================

\titleformat{\chapter}[display]
    {\normalfont\huge\bfseries\color{darkgray}}
    {\chaptertitlename\ \thechapter}
    {20pt}
    {\Huge}

\titleformat{\section}
    {\normalfont\Large\bfseries\color{darkgray}}
    {\thesection}
    {1em}
    {}

\titleformat{\subsection}
    {\normalfont\large\bfseries\color{midgray}}
    {\thesubsection}
    {1em}
    {}

% ============================================================
% BOXEN
% ============================================================

\tcbset{
    colback=boxbg,
    colframe=accentcolor,
    boxrule=0.5pt,
    arc=2pt,
    left=10pt,
    right=10pt,
    top=8pt,
    bottom=8pt
}

\newtcolorbox{szenenbox}[1][]{
    colback=white,
    colframe=darkgray,
    fonttitle=\bfseries,
    title=#1,
    boxrule=0.8pt
}

\newtcolorbox{dialogbox}{
    colback=white,
    colframe=lightgray,
    boxrule=0.3pt,
    left=15pt,
    right=15pt
}

\newtcolorbox{voiceoverbox}{
    colback=boxbg,
    colframe=accentcolor,
    boxrule=0.5pt,
    fontupper=\itshape
}

% ============================================================
% SZENEN-BEFEHLE
% ============================================================

\newcommand{\szene}[2]{%
    \vspace{1em}
    \noindent\textbf{\textsf{#1}}\\
    \textit{#2}
    \vspace{0.5em}
}

\newcommand{\figur}[1]{%
    \vspace{0.5em}
    \noindent\hspace{3cm}\textbf{\MakeUppercase{#1}}\\
    \hspace{3cm}
}

\newcommand{\regieanweisung}[1]{%
    \textit{(#1)}
}

\newcommand{\voiceover}[1]{%
    \begin{voiceoverbox}
    \textbf{JOHANNES (V.O.)}\\
    #1
    \end{voiceoverbox}
}

% ============================================================
% DOKUMENT
% ============================================================

\begin{document}

% ============================================================
% TITELSEITE
% ============================================================

\begin{titlepage}
    \centering
    \vspace*{3cm}

    {\Huge\bfseries\color{darkgray} DIE VERZAUBERUNG}

    \vspace{1cm}

    {\Large Ein Film nach dem Roman von}

    \vspace{0.5cm}

    {\LARGE\textsc{Hermann Broch}}

    \vspace{3cm}

    {\large Drehbuch von}

    \vspace{0.5cm}

    {\Large\bfseries Boris Marte}

    \vfill

    \begin{tabular}{ll}
        Format: & Spielfilm, 120 Minuten \\
        Genre: & Psychologisches Drama \\
        Epoche: & 1930er Jahre, Alpenregion \\
        Fassung: & Januar 2026
    \end{tabular}

    \vspace{2cm}

    {\small\textit{„Anatomie eines kollektiven Wahns"}}

\end{titlepage}

% ============================================================
% IMPRESSUM
% ============================================================

\thispagestyle{empty}
\vspace*{\fill}

\begin{center}
\textsc{Die Verzauberung}\\[0.5em]
Drehbuch nach dem Roman von Hermann Broch\\[2em]

Drehbuch: Boris Marte\\[0.5em]
Fassung: Januar 2026\\[2em]

\textit{Alle Rechte vorbehalten.}\\
\textit{Basierend auf dem Roman „Die Verzauberung" (1935) von Hermann Broch.}
\end{center}

\vspace*{\fill}
\newpage

% ============================================================
% INHALTSVERZEICHNIS
% ============================================================

\tableofcontents
\newpage

% ============================================================
% TEIL I: TREATMENT
% ============================================================

\part{Treatment}

\chapter{Übersicht}

\section{Technische Daten}

\begin{tabularx}{\textwidth}{lX}
\toprule
\textbf{Element} & \textbf{Spezifikation} \\
\midrule
Format & Spielfilm \\
Länge & ca. 120 Minuten \\
Epoche & 1930er Jahre, Alpenregion \\
Erzählform & Voice-Over (Landarzt), chronologisch \\
Tonalität & Psychologischer Realismus \\
Referenzen & Haneke, Dreyer, früher Bergman \\
\bottomrule
\end{tabularx}

\section{Logline}

\begin{tcolorbox}
Ein Landarzt beobachtet machtlos, wie ein charismatischer Fremder sein abgelegenes Bergdorf durch mythische Rhetorik in einen kollektiven Wahn treibt – eine Anatomie der Verführbarkeit in Zeiten des Sinnverlusts.
\end{tcolorbox}

\section{Prämisse}

\textbf{Kernfrage:} Wie werden vernünftige Menschen zu Mittätern?

\textbf{These:} Wenn traditionelle Sinnstrukturen (Religion, Gemeinschaft, Arbeit) ihre Bindungskraft verlieren, entsteht ein Vakuum, das von mythischer Verführung gefüllt wird. Die Vernunft allein ist zu schwach, um dem Rausch der Zugehörigkeit zu widerstehen.

\subsection{Subtexte}
\begin{itemize}
    \item Allegorie auf den aufkommenden Faschismus
    \item Kritik der Massenpsychologie (Le Bon, Freud)
    \item Die Einsamkeit des rationalen Beobachters
    \item Schuld durch Untätigkeit
\end{itemize}

\section{Setting}

\textbf{Ort:} Ein namenloses Bergdorf in den österreichischen oder Schweizer Alpen. 800--1000 Meter Höhe. Etwa 200--300 Einwohner. Erreichbar nur über eine einzige Passstraße, die im Winter oft unpassierbar ist.

\textbf{Zeit:} 1934--1936. Die Weltwirtschaftskrise wirkt nach. In den Städten ist der Faschismus bereits präsent, doch das Dorf lebt noch in einer Zwischenwelt -- nicht mehr traditional, noch nicht modern.

\textbf{Atmosphäre:}
\begin{itemize}
    \item Karge, erhabene Landschaft (keine Postkartenidylle)
    \item Enge Innenräume, niedrige Decken, wenig Licht
    \item Kontrast: weite Bergpanoramen vs. klaustrophobische Dorfstrukturen
    \item Jahreszeiten als dramaturgisches Element: Beginn im Spätsommer, Ende im Winter
\end{itemize}

% ============================================================
% KAPITEL: FIGUREN
% ============================================================

\chapter{Die Figuren}

\section{Dr. Johannes Müller -- Der Landarzt}

\begin{tcolorbox}[title=Basisdaten]
\begin{tabular}{ll}
Alter: & 48 Jahre \\
Herkunft: & Wien, großbürgerlich \\
Beruf: & Landarzt \\
Im Dorf seit: & 15 Jahren \\
Familienstand: & Geschieden, kinderlos (Tochter verstorben)
\end{tabular}
\end{tcolorbox}

\begin{quote}
\textit{„Ich wusste, was geschah. Ich konnte es benennen. Aber Wissen ist keine Macht."}
\end{quote}

Seit 15 Jahren Arzt im Dorf. Ursprünglich aus der Stadt, hat er sich bewusst für die Abgeschiedenheit entschieden -- nach einer gescheiterten Ehe, einer toten Tochter, einem Zusammenbruch. Das Dorf war Flucht und Heilung. Er ist respektiert, aber nicht integriert. Ein Einzelgänger, der liest, denkt, beobachtet.

\textbf{Funktion:} Erzähler, moralisches Zentrum, tragischer Zeuge

\textbf{Charakterbogen:}
\begin{enumerate}
    \item Akt I: Beobachter -- „Ich sehe, aber ich bin nicht beteiligt."
    \item Akt II: Warner -- „Ich sehe die Gefahr und spreche sie aus."
    \item Akt III: Gescheiterter -- „Ich habe gesprochen, aber nicht gehandelt."
\end{enumerate}

\section{Marius Ratti -- Der Wanderer}

\begin{tcolorbox}[title=Basisdaten]
\begin{tabular}{ll}
Alter: & 35--40 Jahre (unklar) \\
Herkunft: & Unbekannt (behauptet: Südtirol) \\
Beruf: & Behauptet: Geologe \\
Besonderheit: & Akzent nicht zuordenbar
\end{tabular}
\end{tcolorbox}

\begin{quote}
\textit{„Ich bringe nichts, was nicht schon in euch ist. Ich wecke nur, was schläft."}
\end{quote}

Niemand weiß, woher er kommt. Er spricht Hochdeutsch mit leichtem, nicht zuordenbarem Akzent. Gepflegtes Äußeres, aber keine Eleganz -- eher eine asketische Strenge. Er behauptet, Geologe zu sein, interessiert an den Bergen. Doch seine wahre Expertise ist das Lesen von Menschen.

\textbf{Funktion:} Antagonist, Katalysator, Verführer

\textbf{Charakteristik:} Nie schreiend, nie gewalttätig, immer ruhig. Seine Gefährlichkeit liegt in der Sanftheit.

\section{Maria Gisson -- Die Sehende}

\begin{tcolorbox}[title=Basisdaten]
\begin{tabular}{ll}
Alter: & 22 Jahre \\
Herkunft: & Geboren im Dorf \\
Familie: & Vater: Wenzel Gisson. Mutter: verstorben \\
Besonderheit: & „Anders" seit Kindheitskrankheit
\end{tabular}
\end{tcolorbox}

\textit{Stumm für den Großteil des Films. Ihr Schweigen ist ihr Widerstand -- und ihr Verhängnis.}

Tochter des reichsten Bauern, aber Außenseiterin. Als Kind hatte sie hohes Fieber, seitdem gilt sie als „anders" -- nicht geisteskrank, aber verträumt, weltentrückt. Der Wanderer erkennt in ihr das perfekte Opfer: rein, anders, wehrlos.

\textbf{Funktion:} Opfer, symbolische Reinheit, das Schweigen der Vernunft

\section{Wenzel Gisson -- Der Bauer}

\begin{quote}
\textit{„Wir haben immer gearbeitet. Aber wofür? Das konnte mir keiner sagen."}
\end{quote}

Marias Vater. Der erfolgreichste Bauer des Dorfes, aber innerlich leer. Seine Frau starb vor Jahren, seine Tochter ist ihm fremd. In Ratti sieht er erstmals jemanden, der Antworten zu haben scheint.

\textbf{Funktion:} Täter, der tragische Vater, Archetyp des Mitläufers

\section{Pfarrer Anton Lindner}

Seit 30 Jahren im Dorf. Müde, aber nicht zynisch. Er glaubt noch, aber sein Glaube erreicht die Menschen nicht mehr.

\textbf{Funktion:} Repräsentant der versagenden Institution

\section{Agathe Meissner -- Die Wirtin}

\begin{quote}
\textit{„Endlich redet einer Klartext. Endlich sagt einer, wo's langgeht."}
\end{quote}

Betreibt das einzige Gasthaus. Witwe, praktisch, handfest. Sie wird zur ersten und glühendsten Anhängerin Rattis.

\textbf{Funktion:} Repräsentantin des „normalen" Mitläufertums

% ============================================================
% KAPITEL: DRAMATURGISCHE ÜBERSICHT
% ============================================================

\chapter{Dramaturgische Übersicht}

\section{Akt I: Die Stille vor dem Sturm}
\textit{Minuten 1--30}

Das Dorf im Spätsommer. Ernte. Oberflächliche Normalität. Der Landarzt führt uns ein (Voice-Over). Wir sehen die Risse: Streit um Weidegrenzen, eine Schlägerei nach dem Kirchgang, nervöse Patienten ohne körperliche Ursache.

\textbf{Schlüsselszene:} Der Arzt behandelt einen Bauern wegen Schlaflosigkeit. Der Mann sagt: „Es ist, als würde etwas kommen. Ich weiß nicht was."

Dann: Ankunft des Fremden. Er kommt zu Fuß über den Pass.

\section{Akt II: Die Verzauberung}
\textit{Minuten 30--90}

Die Versammlungen werden größer. Ratti spricht nun offener: von Reinigung, von Opfer, von Notwendigkeit. Die Sprache wird mythischer. Maria gerät ins Visier.

\textbf{Midpoint (Min. 58--60):} Ratti pflanzt die Idee in Gissons Kopf: Maria ist „ein Zeichen".

\textbf{Schlüsselszene:} Der Arzt versucht, mit Gisson zu sprechen. „Sie sehen doch, wohin das führt." Gisson: „Sie verstehen nicht. Sie sind Wissenschaftler."

\section{Akt III: Das Opfer}
\textit{Minuten 90--120}

Das kranke Kind stirbt. Die Stimmung kippt. Maria wird offen als Ursache benannt. Nächtliche Zeremonie. Fackeln. Der Arzt versucht einzugreifen -- wird zurückgehalten.

\textbf{Marias letzte Worte:} „Du glaubst nicht einmal selbst daran."

\textbf{Rattis Antwort:} „Das ist das Wunder. Es wirkt trotzdem."

Stille. Schnee. Das Dorf nach der Tat. Ratti ist verschwunden.

\textbf{Schlussmonolog:}
\begin{voiceoverbox}
„Ich warte. Nicht auf Erlösung -- die gibt es nicht. Ich warte auf den nächsten Winter. Und auf den nächsten Fremden. Denn die Verzauberung endet nie. Sie schläft nur. Bis jemand kommt, der sie weckt."
\end{voiceoverbox}

% ============================================================
% TEIL II: DREHBUCH
% ============================================================

\part{Drehbuch}

\chapter{Akt I: Die Stille vor dem Sturm}

\section*{Sequenz 1: Das Dorf}

\szene{SZENE 1}{EXT. BERGLANDSCHAFT -- MORGENDÄMMERUNG}

Nebel liegt in den Tälern. Die ersten Sonnenstrahlen berühren die Gipfel. Eine Landschaft von erhabener Gleichgültigkeit.

Die Kamera schwebt langsam nach unten.

Ein Dorf wird sichtbar: kleine Häuser, eine Kirche mit spitzem Turm, Felder in verschiedenen Braun- und Goldtönen. Rauch steigt aus Schornsteinen.

\voiceover{Es gibt Orte, die vergessen wurden. Nicht von der Landkarte -- von der Zeit. Unser Dorf war so ein Ort. Die Welt da unten tobte. Wir wussten das. Aber wir glaubten, die Berge würden uns schützen.}

Die Kamera kommt zum Stillstand.

\textbf{INSERT:} „SEPTEMBER 1934"

\vspace{1em}
\hrule
\vspace{1em}

\szene{SZENE 2}{EXT. FELDER -- MORGEN}

Hände greifen in die Erde. Ein Spaten sticht in den Boden. Kartoffeln werden ausgegraben.

Wir sehen Gesichter: verwittert, verschlossen, müde. Bauern bei der Erntearbeit. Niemand spricht.

Ein Kind (8) treibt Ziegen über einen Pfad. Es singt leise vor sich hin -- ein altes Volkslied, dessen Worte wir nicht ganz verstehen.

\voiceover{Es war Erntezeit. Die letzte gute Zeit, bevor der Winter kam. Bevor alles kam.}

\vspace{1em}
\hrule
\vspace{1em}

\szene{SZENE 3}{INT. ARZTZIMMER -- TAG}

Ein karger Raum: Holzmöbel, medizinische Instrumente, Bücher an den Wänden. An einer Wand ein Anatomie-Poster. Durch das Fenster fällt warmes Herbstlicht.

DR. JOHANNES MÜLLER (48) -- graumeliertes Haar, konzentrierter Blick, sorgfältige Bewegungen -- misst den Blutdruck einer alten BÄUERIN (70).

Die Bäuerin sitzt still. Ihre Augen wandern unruhig durch den Raum.

\begin{dialogbox}
\figur{Johannes}
Der Blutdruck ist gut. Herz und Lunge auch.

\regieanweisung{Er legt das Stethoskop weg.}

\figur{Johannes}
Was führt Sie zu mir?

\figur{Bäuerin}
Ich schlafe nicht mehr, Herr Doktor.

\figur{Johannes}
Seit wann?

\figur{Bäuerin}
Wochen. Ich liege da und... warte.

\figur{Johannes}
Worauf?

\figur{Bäuerin}
\regieanweisung{Ihre Augen sind weit.}
Das ist es ja. Ich weiß es nicht.
\end{dialogbox}

Johannes schreibt ein Rezept. Seine Hand zögert kurz.

\voiceover{Ich verschrieb ihr Baldrian. Es würde nicht helfen. Jakobs Frau war nicht die Einzige. In diesem Herbst kamen sie alle -- mit Schlaflosigkeit, Unruhe, vagen Ängsten. Ich konnte keine Ursache finden. Noch nicht.}

\vspace{1em}
\hrule
\vspace{1em}

\szene{SZENE 10}{EXT. PASSSTRASSE -- NACHMITTAG}

Die einzige Straße, die zum Dorf führt. Sie windet sich den Berg hinauf.

Eine einzelne Gestalt kommt zu Fuß. MARIUS RATTI (38) -- schlanke, aufrechte Haltung. Einfache, aber gepflegte Kleidung. Ein kleiner Rucksack auf dem Rücken.

Sein Gesicht ist schwer zu lesen. Nicht freundlich, nicht unfreundlich. Aufmerksam.

\voiceover{Er kam an einem Donnerstag im September. Zu Fuß, über den Pass. Niemand hatte ihn erwartet. Niemand fragte, woher er kam. Später fragten wir uns, warum nicht.}

Ratti bleibt stehen. Betrachtet das Dorf unter sich. Ein Lächeln -- kaum sichtbar.

Dann geht er weiter.

\vspace{1em}
\hrule
\vspace{1em}

\szene{SZENE 13}{INT. GASTHAUS -- SPÄTER ABEND}

Johannes betritt das Gasthaus. Er trägt seinen Arztkittel noch unter dem Mantel.

Agathe deutet mit dem Kopf auf Ratti.

\begin{dialogbox}
\figur{Agathe}
\regieanweisung{leise}
Der Neue. Interessanter Typ.
\end{dialogbox}

Johannes geht zu Rattis Tisch. Bleibt stehen.

\begin{dialogbox}
\figur{Johannes}
Der neue Gast.

\figur{Ratti}
\regieanweisung{blickt auf}
Ja.

\figur{Johannes}
Ich bin der Arzt hier. Müller.
\regieanweisung{setzt sich}
Darf ich?

\figur{Ratti}
Bitte.
\end{dialogbox}

Sie mustern einander. Zwei intelligente Männer, die den anderen einschätzen.

\begin{dialogbox}
\figur{Ratti}
Marius Ratti. Geologe.

\figur{Johannes}
Geologe. Interessant. Was führt einen Geologen hierher?

\figur{Ratti}
Die Schichtungen. Bemerkenswert, was sich unter der Oberfläche findet.

\figur{Johannes}
Bei Gestein?

\figur{Ratti}
\regieanweisung{Pause, ein Lächeln}
Unter anderem.
\end{dialogbox}

Johannes lehnt sich zurück.

\begin{dialogbox}
\figur{Johannes}
Welches Fachgebiet genau?

\figur{Ratti}
Die Geologie der Seele, wenn Sie so wollen. Was unter der Oberfläche liegt. Bei Gestein. Bei Landschaften.
\regieanweisung{Pause}
Bei Menschen.
\end{dialogbox}

Stille. Die beiden Männer sehen einander an.

\begin{dialogbox}
\figur{Johannes}
Eine ungewöhnliche Spezialisierung.

\figur{Ratti}
Es gibt viel zu entdecken. Gerade hier.

\figur{Johannes}
\regieanweisung{steht auf}
Willkommen im Dorf, Herr Ratti.

\figur{Ratti}
Danke, Herr Doktor. Ich bin sicher, wir werden uns wiedersehen.
\end{dialogbox}

Johannes geht. Ratti sieht ihm nach. Sein Lächeln verblasst nicht -- aber es erreicht seine Augen nicht.

% ============================================================
% AKT II
% ============================================================

\chapter{Akt II: Die Verzauberung}

\section*{Sequenz 5: Die ersten Rituale}

\szene{SZENE 17}{EXT. DORF -- TAG}

Die Bäume haben sich verfärbt. Goldene und rote Blätter. Das Licht ist härter geworden.

\textbf{INSERT:} „OKTOBER 1934"

\voiceover{Es begann langsam. So langsam, dass ich es fast übersehen hätte. Aber die Verzauberung braucht keine Eile.}

\vspace{1em}
\hrule
\vspace{1em}

\szene{SZENE 20}{INT. GASTHAUS -- ABEND (DRITTE VERSAMMLUNG)}

Dreißig Menschen. Der Raum ist voll.

Johannes steht am Rand. Er ist nicht Teil der Gruppe -- aber er kann nicht wegsehen.

\begin{dialogbox}
\figur{Ratti}
Ihr fragt, was ich will? Ich will nichts. Ich bin nur ein Wanderer.

\figur{Stimme}
Aber Sie wissen Dinge.

\figur{Ratti}
Ich weiß nur, was ihr mir erzählt. Und ich sehe, was ihr nicht seht.

\figur{Agathe}
Was sehen Sie?

\figur{Ratti}
\regieanweisung{Pause, sieht alle an}
Eine Gemeinschaft, die vergessen hat, was sie ist. Menschen, die nach Sinn hungern. Ein Dorf, das schläft -- und darauf wartet, geweckt zu werden.

\figur{Gisson}
\regieanweisung{aufstehend}
Dann wecken Sie uns.
\end{dialogbox}

Ratti lächelt. Nicht triumphierend -- fast bescheiden.

\begin{dialogbox}
\figur{Ratti}
Nicht ich. Wir. Zusammen.
\end{dialogbox}

Applaus. Johannes am Rand -- sein Gesicht: Entsetzen, das er verbirgt.

\voiceover{Er hatte nichts gesagt. Keine Fakten, keine Beweise, keine Versprechen. Und doch -- sie waren ihm bereits verfallen.}

\vspace{1em}
\hrule
\vspace{1em}

\szene{SZENE 21}{EXT. WALDRAND -- TAG}

Johannes findet Ratti allein auf einem Felsen sitzend, den Blick auf die Berge gerichtet.

\begin{dialogbox}
\figur{Johannes}
\regieanweisung{ohne Begrüßung}
Was wollen Sie hier, Ratti?

\figur{Ratti}
\regieanweisung{dreht sich nicht um}
Guten Tag, Herr Doktor.

\figur{Johannes}
Ich habe Sie etwas gefragt.
\end{dialogbox}

Ratti steht auf. Wendet sich um. Sein Gesicht: ruhig, fast freundlich.

\begin{dialogbox}
\figur{Ratti}
Dasselbe wie Sie. Verstehen.

\figur{Johannes}
Ich heile Menschen.

\figur{Ratti}
Heilen Sie? Wirklich?
\regieanweisung{tritt näher}
Sie geben Pillen gegen Schlaflosigkeit. Tropfen gegen Unruhe. Aber die Krankheit... die behandeln Sie nicht.

\figur{Johannes}
Und Sie können es?

\figur{Ratti}
Ich versuche es. Mit anderen Mitteln.

\figur{Johannes}
Mit Mythen. Mit Aberglauben.

\figur{Ratti}
\regieanweisung{lächelt}
Mit dem, was wirkt.
\regieanweisung{Pause}
Diese Menschen sind seelenkrank. Sie brauchen nicht Ihre Pillen. Sie brauchen... Bedeutung.
\end{dialogbox}

Stille. Die beiden Männer stehen sich gegenüber.

\begin{dialogbox}
\figur{Ratti}
Sie könnten mitmachen, Herr Doktor. Wir brauchen kluge Köpfe.

\figur{Johannes}
Nein.

\figur{Ratti}
\regieanweisung{nickt}
Schade. Aber ich verstehe. Der Beobachter will Beobachter bleiben.
\regieanweisung{dreht sich weg}
Nur -- irgendwann muss man sich entscheiden. Zuschauen ist auch eine Entscheidung.
\end{dialogbox}

Er geht. Johannes bleibt allein zurück.

\vspace{1em}
\hrule
\vspace{1em}

\szene{SZENE 29}{EXT. BACHUFER -- TAG}

MARIA (22) sitzt am Wasser. Sie wirft Steine. Konzentrische Kreise breiten sich aus.

Johannes nähert sich vorsichtig.

\begin{dialogbox}
\figur{Johannes}
Maria?
\end{dialogbox}

Sie reagiert nicht sofort. Dann dreht sie sich langsam um.

\begin{dialogbox}
\figur{Johannes}
Wie geht es dir?

\figur{Maria}
\regieanweisung{leise}
Wie es mir geht?

\figur{Johannes}
Ich mache mir Sorgen. Du musst weg von hier.
\end{dialogbox}

Stille.

\begin{dialogbox}
\figur{Johannes}
Ich kann dich in die Stadt bringen. Zu Verwandten. Dort bist du sicher.

\figur{Maria}
\regieanweisung{schüttelt den Kopf}
Die Stadt?

\figur{Johannes}
Hier ist es nicht sicher. Nicht für dich.

\figur{Maria}
\regieanweisung{sieht zum Wasser}
In der Stadt ist es auch.

\figur{Johannes}
Was ist auch?
\end{dialogbox}

Lange Pause. Maria spricht zum ersten Mal mehr als ein paar Worte.

\begin{dialogbox}
\figur{Maria}
Die Angst. Und was aus ihr wächst.
\end{dialogbox}

Johannes ist erschüttert.

\begin{dialogbox}
\figur{Johannes}
Was siehst du, Maria?

\figur{Maria}
\regieanweisung{dreht sich weg}
Zu viel.
\end{dialogbox}

Sie steht auf und geht. Johannes bleibt am Ufer.

\voiceover{Sie hatte mehr gesagt als in Jahren zuvor. Und doch -- ich verstand erst später, was sie meinte.}

% ============================================================
% AKT III
% ============================================================

\chapter{Akt III: Das Opfer}

\section*{Sequenz 9: Die Zeremonie}

\szene{SZENE 37}{EXT. BERGPLATEAU -- NACHT}

Fackeln bilden einen Kreis. Das alte Steinkreuz ragt auf. Dreißig, vierzig Menschen. In der Mitte: Maria in einem weißen Kleid.

Ratti tritt vor.

\begin{dialogbox}
\figur{Ratti}
Wir sind hier, weil wir reinigen müssen. Nicht mit Hass. Mit Liebe. Die Reine geht voran. Die Reine öffnet den Weg.
\end{dialogbox}

Maria sieht ihn an. Kein Hass. Keine Angst. Nur: Verstehen.

\begin{dialogbox}
\figur{Maria}
\regieanweisung{leise, aber deutlich}
Du glaubst nicht einmal selbst daran.
\end{dialogbox}

Ratti erstarrt. Sein Lächeln verrutscht -- für einen Moment.

Dann fängt er sich.

\begin{dialogbox}
\figur{Ratti}
\regieanweisung{nur für sie}
Das ist das Wunder. Es wirkt trotzdem.
\end{dialogbox}

Maria nickt langsam.

\voiceover{Das waren ihre letzten Worte. Sie war die Einzige, die ihn durchschaut hatte.}

Johannes reißt sich los. Stürzt nach vorne.

\begin{dialogbox}
\figur{Johannes}
Nein! Hört auf!
\end{dialogbox}

Er wird zu Boden geworfen. Gesichter über ihm.

\begin{dialogbox}
\figur{Gisson}
\regieanweisung{steht über ihm}
Schweigen Sie. Bitte. Machen Sie es nicht schwerer.
\end{dialogbox}

Johannes sieht in Gissons Augen: Der Mann ist nicht mehr erreichbar.

Was dann geschieht -- wir sehen es nicht. Nur:

\begin{itemize}
    \item Gissons Gesicht, der wegsieht.
    \item Agathes Lippen, die sich im Gebet bewegen.
    \item Der Pfarrer am Rand, erstarrt.
    \item Schatten auf dem Steinkreuz.
\end{itemize}

Dann: Stille.

Langsame \textbf{SCHWARZBLENDE}.

\textbf{HALTEN.}

Dann: Vogelgezwitscher.

\vspace{1em}
\hrule
\vspace{1em}

\section*{Sequenz 11: Epilog}

\szene{SZENE 44 (SCHLUSS)}{INT. ARZTZIMMER -- NACHT}

Johannes steht am Fenster. Draußen: Schnee, Stille, Dunkelheit.

Ein Licht in der Ferne -- Gissons Hof.

Dann: Bewegung auf der Straße. Eine Gestalt.

Johannes beobachtet. Wird aufmerksam.

Die Gestalt kommt näher. Ein Mann? Eine Frau? Wir können es nicht erkennen.

Ein Fremder? Oder nur ein Dorfbewohner, der spät nach Hause geht?

Johannes' Gesicht: wachsam. Nicht ängstlich -- wissend.

\voiceover{Ich warte. Nicht auf Erlösung -- die gibt es nicht. Ich warte auf den nächsten Winter. Und auf den nächsten Fremden.}

Die Gestalt verschwindet hinter einem Haus.

Johannes bleibt am Fenster.

\voiceover{Denn die Verzauberung endet nie. Sie schläft nur.\\[0.5em]
Bis jemand kommt, der sie weckt.}

Sein Gesicht. Die Dunkelheit draußen. Der Schnee.

Langsame \textbf{ABBLENDE}.

\textbf{SCHWARZ.}

\begin{center}
\vspace{2em}
{\Large\textbf{ENDE}}
\vspace{2em}
\end{center}

% ============================================================
% TEIL III: PRODUKTIONSINFORMATIONEN
% ============================================================

\part{Produktionsinformationen}

\chapter{Visuelles Konzept}

\section{Leitprinzipien}

Der Film folgt einer visuellen Strategie, die die Themen der Erzählung spiegelt: \textbf{Beobachtung, Einschließung, Entmenschlichung}.

\begin{tabularx}{\textwidth}{lX}
\toprule
\textbf{Prinzip} & \textbf{Umsetzung} \\
\midrule
Distanz & Lange Brennweiten, Beobachterperspektive \\
Statik & Wenig Kamerabewegung, lange Einstellungen \\
Natürliches Licht & Keine künstliche Ausleuchtung \\
Reduktion & Farbpalette wird von Akt zu Akt entsättigter \\
Gesichter & Nahaufnahmen als moralische Landschaften \\
\bottomrule
\end{tabularx}

\section{Farbdramaturgie}

\begin{description}
    \item[Akt I:] Warme Herbsttöne: Gold, Braun, gedämpftes Grün
    \item[Akt II:] Zunehmend entsättigt, Braun dominiert
    \item[Akt III:] Fast monochrom, Blau-Grau (Schnee, Winter)
\end{description}

\section{Tonale Referenzen}

\begin{tabularx}{\textwidth}{lX}
\toprule
\textbf{Film} & \textbf{Element} \\
\midrule
Das weiße Band (Haneke) & Dorfmilieu, versteckte Gewalt, Kälte \\
Ordet (Dreyer) & Religiöse Spannung, lange Einstellungen \\
Das Fest (Vinterberg) & Kollektive Schuld, Familienstruktur \\
The Witch (Eggers) & Isoliertes Setting, mythische Atmosphäre \\
Heimat (Reitz) & Authentische 1930er-Atmosphäre \\
\bottomrule
\end{tabularx}

\chapter{Budget-Übersicht}

\section{Gesamtbudget}

\begin{tcolorbox}
\centering
{\Large\textbf{€ 4.200.000 -- € 5.100.000}}\\[0.5em]
Koproduktion: Deutschland / Österreich / Schweiz
\end{tcolorbox}

\section{Aufschlüsselung}

\begin{tabularx}{\textwidth}{Xlrr}
\toprule
\textbf{Kategorie} & \textbf{Min (€)} & \textbf{Max (€)} & \textbf{\%} \\
\midrule
Vorproduktion & 180.000 & 220.000 & 4\% \\
Drehstab & 620.000 & 750.000 & 15\% \\
Darsteller & 480.000 & 600.000 & 12\% \\
Ausstattung & 380.000 & 480.000 & 10\% \\
Technik & 420.000 & 520.000 & 10\% \\
Locations & 240.000 & 320.000 & 6\% \\
Reise/Unterkunft & 320.000 & 400.000 & 8\% \\
Postproduktion & 480.000 & 580.000 & 12\% \\
Musik & 120.000 & 160.000 & 3\% \\
Versicherung/Recht & 140.000 & 180.000 & 4\% \\
Allgemein & 240.000 & 300.000 & 6\% \\
Unvorhergesehenes & 400.000 & 370.000 & 6\% \\
\midrule
\textbf{GESAMT} & \textbf{4.200.000} & \textbf{5.100.000} & \textbf{100\%} \\
\bottomrule
\end{tabularx}

\chapter{Szenenindex}

\begin{longtable}{clllr}
\toprule
\textbf{Nr.} & \textbf{Ort} & \textbf{Zeit} & \textbf{Inhalt} & \textbf{Min.} \\
\midrule
\endhead
1 & Berglandschaft & Morgen & Eröffnung, Dorf von oben & 1:30 \\
2 & Felder & Morgen & Erntearbeit & 1:00 \\
3 & Arztzimmer & Tag & Schlaflose Bäuerin & 2:30 \\
4 & Kirchplatz & Tag & Streit nach Messe & 1:30 \\
5 & Bergwiese & Tag & Maria allein & 2:00 \\
10 & Passstraße & Nachmittag & Rattis Ankunft & 1:00 \\
13 & Gasthaus & Abend & Erste Begegnung J./R. & 3:00 \\
20 & Gasthaus & Abend & Dritte Versammlung & 3:00 \\
21 & Waldrand & Tag & Konfrontation J./R. & 4:00 \\
23 & Waldlichtung & Nacht & Erstes Ritual & 4:00 \\
26 & Gisson-Hof & Abend & R. und Gisson & 4:00 \\
29 & Bachufer & Tag & Maria spricht & 3:30 \\
34 & Rattis Zimmer & Nacht & Letzte Konfrontation & 4:00 \\
37 & Bergplateau & Nacht & Das Opfer & 6:00 \\
41 & Gisson-Hof & Winter & Gespräch mit Gisson & 3:00 \\
44 & Arztzimmer & Nacht & Schlussbild & 3:00 \\
\midrule
& & & \textbf{GESAMT} & \textbf{~120:00} \\
\bottomrule
\end{longtable}

% ============================================================
% ANHANG
% ============================================================

\appendix

\chapter{Voice-Over Übersicht}

\begin{longtable}{cl}
\toprule
\textbf{Szene} & \textbf{Thema} \\
\midrule
\endhead
1 & Einführung Dorf \\
3 & Einführung Krankheit \\
5 & Einführung Maria \\
10 & Einführung Ratti \\
14 & Rattis Methode \\
20 & Die Verführung \\
23 & Das erste Ritual \\
29 & Marias Warnung \\
37 & Die Schuld \\
44 & Der Schlussmonolog \\
\bottomrule
\end{longtable}

\chapter{Kontakt}

\begin{tcolorbox}
\begin{center}
\textbf{DIE VERZAUBERUNG}\\
Ein Film nach dem Roman von Hermann Broch\\[1em]

\textbf{Drehbuch:} Boris Marte\\[1em]

\textit{[Produktionsfirma]}\\
\textit{[Adresse]}\\
\textit{[E-Mail]}\\
\textit{[Telefon]}
\end{center}
\end{tcolorbox}

\vfill

\begin{center}
\textit{„Die Verzauberung endet nie. Sie schläft nur."}
\end{center}

\end{document}
